\documentclass[11pt]{article}
\usepackage[margin=1.1in]{geometry}
\usepackage{amsmath, amssymb, amsthm}
\usepackage{mathtools}
\usepackage{microtype}
\usepackage{parskip}
\usepackage{bm}
\usepackage[hidelinks]{hyperref}

\newtheorem{theorem}{Theorem}
\newtheorem{lemma}[theorem]{Lemma}
\newtheorem{proposition}[theorem]{Proposition}
\theoremstyle{remark}
\newtheorem*{remark}{Remark}


\title{Task-Background Document:\\ From the Free-Fermion Action to the $G$--$\Sigma$ and\\ Coadjoint-Orbit Bosonization Actions}
\author{}
\date{}

\begin{document}
\maketitle

\noindent
This document fixes the starting point and the two intermediate targets used
throughout the project. Section~\ref{sec:gsigma} derives the bilocal
($G$--$\Sigma$) action of free, gapless, nonrelativistic fermions from the
microscopic coherent-state action, and identifies its saddle point. Section~\ref{sec:coadjoint}
derives the coadjoint-orbit bosonization action from the same microscopic action,
following the Appendix of Ye and Wang, \emph{Phys.\ Rev.\ B} \textbf{112}, 035113
(2025). The conventions in \S\ref{sec:prelim} are common to both.

\section{From the free-fermion action to the \texorpdfstring{$G$--$\Sigma$}{G-Sigma} action}
\label{sec:gsigma}

\subsection{Preliminaries and conventions}
\label{sec:prelim}

We work with spinless nonrelativistic fermions in $d$ spatial dimensions at
inverse temperature $\beta$, in the imaginary-time coherent-state path integral.
The dynamical variables are Grassmann fields $\psi(\bm{x},\tau)$, $\bar{\psi}(\bm{x},\tau)$
that are antiperiodic in imaginary time, $\psi(\bm{x},\tau+\beta)=-\psi(\bm{x},\tau)$.

To keep the bilocal structure visible we use a collective coordinate
$1 \equiv (\bm{x}_1,\tau_1)$, with
\begin{equation}
\int \mathrm{d} 1 \;\equiv\; \int_0^\beta \mathrm{d}\tau_1 \int \mathrm{d}^d x_1 ,
\qquad
\delta(1,2) \;\equiv\; \delta(\tau_1-\tau_2)\,\delta^d(\bm{x}_1-\bm{x}_2).
\end{equation}
The single-particle Hamiltonian is $\hat h = \varepsilon(-i\bm{\nabla})-\mu$, with
band dispersion $\varepsilon(\bm{k})$, chemical potential $\mu$, and
$\xi_{\bm{k}}\coloneqq \varepsilon(\bm{k})-\mu$. ``Gapless'' means the Fermi surface
$\{\bm{k}:\xi_{\bm{k}}=0\}$ is nonempty and codimension one. The microscopic action is
\begin{equation}
\label{eq:S0}
S_0[\bar{\psi},\psi]
= \int \mathrm{d} 1\;\bar{\psi}(1)\,\big(\partial_{\tau_1}+\hat h_1\big)\,\psi(1).
\end{equation}

It is convenient to write \eqref{eq:S0} as a bilinear $S_0=-\bar{\psi}\,G_0^{-1}\,\psi$
in the free inverse propagator, defined as the operator with kernel
\begin{equation}
\label{eq:G0inv}
G_0^{-1}(1,2) \;\coloneqq\; -\big(\partial_{\tau_1}+\hat h_1\big)\,\delta(1,2).
\end{equation}
The sign in \eqref{eq:G0inv} is chosen so that the Matsubara propagator takes
its standard form: with $\partial_\tau \mapsto -i\omega_n$ on antiperiodic
frequencies $\omega_n=(2n+1)\pi/\beta$,
\begin{equation}
\label{eq:G0matsubara}
G_0(\bm{k},i\omega_n) = \frac{1}{i\omega_n-\xi_{\bm{k}}}.
\end{equation}
We define traces of bilocal kernels by
$\operatorname{Tr}(AB)\coloneqq \int \mathrm{d}1\,\mathrm{d}2\, A(1,2)\,B(2,1)$, and $\operatorname{Tr}\ln$ through its
series. Throughout, $G_0$ denotes the free \emph{propagator}
\eqref{eq:G0matsubara} and $G_0^{-1}$ the operator \eqref{eq:G0inv}; note
$G_0^{-1}=-(\partial_\tau+\hat h)$ carries an explicit minus sign relative to
$\partial_\tau+\hat h$.

\subsection{The bilocal field, and why to introduce it}

The low-energy excitations of a Fermi sea are neutral particle--hole pairs, not
single fermions. Their natural carrier is the equal-time, and more generally the
bilocal, two-point function $\bar{\psi}(2)\psi(1)$: smooth deformations of the
occupied region of phase space are encoded in its fluctuations. We therefore
promote the bilocal to an independent field rather than treating it as a
composite of $\psi,\bar{\psi}$.

Concretely, introduce a bosonic bilocal field $G(1,2)$ together with the
constraint that it equal the microscopic bilinear,
\begin{equation}
\label{eq:Gdef}
G(1,2) \;=\; -\,\psi(1)\,\bar{\psi}(2),
\end{equation}
where the minus sign is dictated by \eqref{eq:G0matsubara}: at the level of
expectation values \eqref{eq:Gdef} gives
$\langle G(1,2)\rangle=-\langle\psi(1)\bar{\psi}(2)\rangle=-\langle T_\tau\,\psi(1)\bar{\psi}(2)\rangle=G_0(1,2)$,
the conventional Green's function. For an \emph{interacting} system this step is
forced: a quartic interaction is decoupled by a field conjugate to a fermion
bilinear, and $G(1,2)$ is the most general such collective field (it carries
every particle--hole channel, not only the density). For the free system there
is nothing to decouple, but the move still buys the essential thing --- it trades
the Grassmann fields for a bosonic bilocal whose path integral is \emph{exact}
and whose fluctuation spectrum is precisely the particle--hole continuum. That
is the arena \S\ref{sec:coadjoint} reorganizes into geometric variables.

\subsection{The self-energy as a Lagrange multiplier}

We enforce \eqref{eq:Gdef} inside the path integral by a functional
delta-function, represented through a second bilocal field $\Sigma(1,2)$ --- the
self-energy --- integrated along the imaginary direction:
\begin{equation}
\label{eq:resolution}
1 = \int \mathcal{D}G\,\mathcal{D}\Sigma\;
\exp\!\Big(-\!\int \mathrm{d}1\,\mathrm{d}2\;\Sigma(2,1)\big[\,G(1,2)+\psi(1)\bar{\psi}(2)\,\big]\Big).
\end{equation}
The $\Sigma$-integral produces $\prod_{1,2}\delta\!\big(G(1,2)+\psi(1)\bar{\psi}(2)\big)$,
which collapses to the constraint \eqref{eq:Gdef}; $\Sigma$ is the field conjugate
to the bilocal, and will turn out to be the self-energy. Inserting
\eqref{eq:resolution} into $Z=\int\mathcal{D}\bar{\psi}\,\mathcal{D}\psi\,e^{-S_0}$
leaves the action quadratic in the fermions.

\subsection{Integrating out the fermions}

Collecting the fermion-bilinear terms from $-S_0$ and from
\eqref{eq:resolution}, and using $\psi(1)\bar{\psi}(2)=-\bar{\psi}(2)\psi(1)$,
\begin{equation}
-S_0 - \int \Sigma(2,1)\,\psi(1)\bar{\psi}(2)
= \bar{\psi}\,G_0^{-1}\,\psi + \bar{\psi}\,\Sigma\,\psi
= \bar{\psi}\,\big(G_0^{-1}+\Sigma\big)\,\psi,
\end{equation}
with $\Sigma$ read as the operator of kernel $\Sigma(2,1)$. The Grassmann
Gaussian integral $\int \mathcal{D}\bar{\psi}\,\mathcal{D}\psi\,
\exp[\bar{\psi}(G_0^{-1}+\Sigma)\psi] = \det\!\big(G_0^{-1}+\Sigma\big)
= e^{\operatorname{Tr}\ln(G_0^{-1}+\Sigma)}$ (an additive, $G,\Sigma$-independent constant is
dropped) leaves a purely bosonic theory,
\begin{equation}
Z = \int \mathcal{D}G\,\mathcal{D}\Sigma\;
\exp\!\big(\operatorname{Tr}\ln(G_0^{-1}+\Sigma) - \operatorname{Tr}(\Sigma G)\big),
\qquad
\operatorname{Tr}(\Sigma G)\coloneqq \int \mathrm{d}1\,\mathrm{d}2\;\Sigma(2,1)\,G(1,2).
\end{equation}

\subsection{The \texorpdfstring{$G$--$\Sigma$}{G-Sigma} action and its saddle point}

We have arrived at the exact bilocal action
\begin{equation}
\label{eq:gsigma}
\boxed{\;
S[G,\Sigma] = -\operatorname{Tr}\ln\!\big(G_0^{-1}+\Sigma\big) + \operatorname{Tr}(\Sigma G) + S_{\mathrm{int}}[G]\;}
\end{equation}
where we have restored $S_{\mathrm{int}}[G]$, the interaction rewritten in terms
of the bilocal (for a density--density interaction $\tfrac12\!\int V\,\rho\rho$
this is a local functional of $G$); for the free fermions of interest
$S_{\mathrm{int}}=0$. Equation \eqref{eq:gsigma} is the central object of this
section: it is a \emph{bosonic} action, obtained without approximation, whose
two fields are the propagator $G$ and the self-energy $\Sigma$.

The saddle-point equations follow from $\delta S/\delta\Sigma=0$ and
$\delta S/\delta G=0$. Using $\delta\,\operatorname{Tr}\ln(G_0^{-1}+\Sigma)=\operatorname{Tr}[(G_0^{-1}+\Sigma)^{-1}\delta\Sigma]$,
\begin{align}
\frac{\delta S}{\delta \Sigma(2,1)} &= -\,(G_0^{-1}+\Sigma)^{-1}(1,2) + G(1,2) = 0
&&\Longrightarrow&& G = \big(G_0^{-1}+\Sigma\big)^{-1},
\label{eq:dyson}\\[2pt]
\frac{\delta S}{\delta G(1,2)} &= \Sigma(2,1) + \frac{\delta S_{\mathrm{int}}}{\delta G(1,2)} = 0
&&\Longrightarrow&& \Sigma = -\,\frac{\delta S_{\mathrm{int}}}{\delta G}.
\label{eq:selfenergy}
\end{align}
Equation \eqref{eq:dyson} is the Dyson equation $G^{-1}=G_0^{-1}+\Sigma$, and
\eqref{eq:selfenergy} identifies $\Sigma$ as the self-energy generated by the
interaction. For the free theory $S_{\mathrm{int}}=0$, so
\begin{equation}
\label{eq:freesaddle}
\Sigma_\star = 0, \qquad G_\star = G_0,
\end{equation}
the free propagator \eqref{eq:G0matsubara}: the saddle of \eqref{eq:gsigma}
exactly reproduces the microscopic two-point function, as it must.

\begin{remark}
The Dyson equation appears here as $G^{-1}=G_0^{-1}+\Sigma$ rather than the
textbook $G^{-1}=G_0^{-1}-\Sigma$; the two differ only by the naming convention
$\Sigma_{\text{here}}=-\Sigma_{\text{textbook}}$, fixed by the sign of the
multiplier term in \eqref{eq:resolution}. Likewise the sign of $\operatorname{Tr}(\Sigma G)$ is
correlated with the minus in the bilocal definition \eqref{eq:Gdef}. The physical
content --- the Dyson equation and the saddle \eqref{eq:freesaddle} --- is
convention-independent.
\end{remark}

\subsection{What the fluctuations are, and the bridge to \texorpdfstring{\S\ref{sec:coadjoint}}{Section 2}}
\label{sec:bridge}

The free saddle \eqref{eq:freesaddle} is trivial; the content is in the
\emph{fluctuations} of $G$ about $G_0$, which \eqref{eq:gsigma} controls exactly.
Writing $G=G_0+\delta G$ and eliminating $\Sigma$ through \eqref{eq:dyson}, the
Gaussian fluctuations of $\delta G$ are the particle--hole pairs of the Fermi
sea; their propagator is the free density--density response, and resumming them
is the bilocal route to the random-phase approximation.

The nonperturbative reorganization used in \S\ref{sec:coadjoint} rests on a
structural fact about the saddle. At zero temperature and equal times the free
propagator reduces to the occupation $n_{\bm{k}}\in\{0,1\}$, so as a single-particle
operator $\hat n$ it is a \emph{projector}, $\hat n^2=\hat n$. The configurations
of $G$ that the dynamics explores at low energy are not arbitrary: they are the
images of $\hat n$ under canonical transformations of single-particle phase space
(area-preserving deformations of the occupied region). That set is an orbit ---
the coadjoint orbit of the group of canonical transformations through the
Fermi-sea projector --- and restricting the $G$--$\Sigma$ action to it, then
expanding $-\operatorname{Tr}\ln$ in gradients, produces the Berry-phase (Kirillov--Kostant--Souriau)
and Hamiltonian terms of the bosonized action. This is the Fermi-liquid analogue
of the $G$--$\Sigma\to$ Schwarzian reduction in SYK, with the coadjoint orbit of
canonical transformations in place of $\mathrm{Diff}(S^1)$. Section~\ref{sec:coadjoint}
carries this out.

\section{From the free-fermion action to the coadjoint-orbit bosonization action}
\label{sec:coadjoint}

This section transcribes the derivation in the Appendix of Ye and Wang,
\emph{Phys.\ Rev.\ B} \textbf{112}, 035113 (2025), specialized to zero magnetic
field. The aim is the same as \S\ref{sec:gsigma} --- recast the free-fermion
theory in collective bosonic variables --- but the route differs.
Section~\ref{sec:gsigma} used the Grassmann coherent state and produced the
bilocal $G$--$\Sigma$ action. Here we use the \emph{coherent state of the Fermi
surface}: a group coherent state built from particle--hole deformations of the
filled sea, the many-body analogue of a spin coherent state. The result is a
manifestly geometric action whose target manifold is the coadjoint orbit of
canonical transformations. Connecting this construction to the $G$--$\Sigma$
action of \S\ref{sec:gsigma} is the project's task; this document establishes the
two endpoints independently.

\subsection{Phase space, star product, and Moyal bracket}

Single-particle observables of a $d$-dimensional system are functions
$F(\bm{x},\bm{p})$ on phase space (the paper works in $d=2$). The reason phase space is
the right arena is that the collective variable will be the occupied region of
$(\bm{x},\bm{p})$, and its deformations are canonical transformations. Operator
multiplication maps not to pointwise multiplication of phase-space symbols but to
the \emph{star product}
\begin{equation}
\label{eq:star}
F\star G = F\,\exp\!\Big(\tfrac{i\hbar}{2}\big(\overleftarrow{\partial}_{\bm{x}}\overrightarrow{\partial}_{\bm{p}}-\overleftarrow{\partial}_{\bm{p}}\overrightarrow{\partial}_{\bm{x}}\big)\Big)\,G ,
\end{equation}
and operator commutators map to the \emph{Moyal bracket}
\begin{equation}
\label{eq:moyal}
\{F,G\}_{\mathrm{M}} = -i\big(F\star G - G\star F\big)
= 2F\,\sin\!\Big(\tfrac{\hbar}{2}\big(\overleftarrow{\partial}_{\bm{x}}\overrightarrow{\partial}_{\bm{p}}-\overleftarrow{\partial}_{\bm{p}}\overrightarrow{\partial}_{\bm{x}}\big)\Big)G .
\end{equation}
Functions $F$ equipped with $\{\cdot,\cdot\}_{\mathrm{M}}$ form the Moyal algebra
$\mathfrak g$, the Lie algebra of canonical transformations; distribution
functions $f$ live in the dual space $\mathfrak g^*$, paired with observables
through
\begin{equation}
\label{eq:pairing}
\langle f, F\rangle \coloneqq \int \frac{\mathrm{d}^dx\,\mathrm{d}^dp}{(2\pi)^d}\,
f(\bm{x},\bm{p})\,F(\bm{x},\bm{p}).
\end{equation}
In the semiclassical regime $\hbar\,\bm{\nabla}_{\bm{x}}\!\cdot\!\bm{\nabla}_{\bm{p}}\ll 1$ the star
product reduces to ordinary multiplication and the Moyal bracket to the Poisson
bracket $\{F,G\}=\partial_{\bm{x}}F\!\cdot\!\partial_{\bm{p}}G-\partial_{\bm{p}}F\!\cdot\!\partial_{\bm{x}}G$.
Keeping the full Moyal algebra (not its Poisson truncation) is what later
distinguishes the quantum theory; we set $\hbar=1$ from here on except where shown.

\subsection{The coadjoint orbit and the ground-state distribution}

Canonical transformations form a group $\mathcal G$ whose Lie algebra is the Moyal
algebra. Generated by $U=\exp(i\phi)$ with $\phi\in\mathfrak g$, they act on a
distribution by conjugation,
\begin{equation}
\label{eq:orbit}
f(\bm{x},\bm{p}) = U \star f_0 \star U^{-1} \;\equiv\; \mathrm{Ad}^*_U\, f_0 ,
\end{equation}
the coadjoint action on $\mathfrak g^*$. The base point is the ground-state
distribution --- the filled Fermi sea,
\begin{equation}
\label{eq:f0}
f_0(\bm{x},\bm{p}) = \Theta\!\big(p_F - |\bm{p}|\big),
\end{equation}
which is the phase-space (Wigner) image of the equal-time Fermi-sea projector
$\hat n$ of \S\ref{sec:bridge}. The set of all distributions reachable from $f_0$,
$\mathcal O_{f_0}=\{\mathrm{Ad}^*_U f_0 : U\in\mathcal G\}$, is the
\emph{coadjoint orbit}; it is the target manifold of the effective theory. The
physical content of \eqref{eq:orbit} is that low-energy configurations are exactly
the area-preserving deformations of the occupied region --- no more, no less.

\subsection{The target action}

The bosonized action of the noninteracting Fermi surface is the sum of a
Berry-phase term and a Hamiltonian term,
\begin{equation}
\label{eq:Saction}
S[f] = S_{\mathrm{WZW}}[f] + S_H[f],\qquad
S_{\mathrm{WZW}} = \int\!\mathrm{d} t\,\big\langle f_0,\, U^{-1}\!\star i\partial_t U\big\rangle,
\qquad
S_H = -\int\!\mathrm{d} t\,\big\langle f,\, \epsilon(\bm{p})\big\rangle,
\end{equation}
with $\epsilon(\bm{p})=\varepsilon_0(\bm{p})-\mu$ the dispersion measured from the
chemical potential (the $\xi_{\bm{k}}$ of \S\ref{sec:prelim}). The WZW term is the
path integral of the Kirillov--Kostant--Souriau symplectic form on the orbit ---
the Berry phase accumulated as the occupied region deforms in time --- and $S_H$
is the single-particle energy. Its equation of motion $\delta S/\delta f=0$ is the
collisionless quantum Boltzmann equation $\partial_t f = \{\epsilon,f\}_{\mathrm{M}}$. The
remainder of this section derives \eqref{eq:Saction} from the microscopic theory.

\subsection{Coherent state of the Fermi surface and its path integral}

Let $|\mathrm{FS}\rangle$ denote the filled Fermi sea. A coherent state is
obtained by dressing it with particle--hole deformations generated by fermion
bilinears,
\begin{equation}
\label{eq:cohstate}
|\mathcal U\rangle = \hat{\mathcal U}\,|\mathrm{FS}\rangle,
\qquad
\hat{\mathcal U} = \exp\!\big(i\phi_{ij}\,c_i^\dagger c_j\big)
\quad (i,j\in[1,L^d]),
\end{equation}
the many-body analogue of a spin coherent state. Via the algebra of bilinears the
$\hat{\mathcal U}$ form a Lie group, and free-fermion evolution
$e^{i\hat H\,\mathrm{d} t}$ commutes with the group average; Schur's lemma then gives a
resolution of identity over the Haar measure $\mathrm{d}\mathcal U$ (up to a constant),
\begin{equation}
\label{eq:resolution2}
\int \mathrm{d}\mathcal U\, |\mathcal U\rangle\langle\mathcal U| = 1 .
\end{equation}
Inserting \eqref{eq:resolution2} at every infinitesimal time slice --- exactly as
for the spin path integral --- yields
\begin{equation}
\label{eq:Zpath}
Z = \int \mathcal D\mathcal U(t)\,
\exp\!\Big[\int \mathrm{d} t\;
\big\langle \mathrm{FS}\big|\,\hat{\mathcal U}^{-1}
\big(-\partial_t - i\hat H\big)\hat{\mathcal U}\,\big|\mathrm{FS}\big\rangle\Big].
\end{equation}

\subsection{Reduction to first-quantized operators}

The exponent of \eqref{eq:Zpath} is a sum of nested commutators of bilinears.
These close on themselves: for $\hat A=c_i^\dagger a_{ij}c_j$ and
$\hat B=c_i^\dagger b_{ij}c_j$, fermion statistics give
\begin{equation}
\label{eq:bilinearcomm}
[\hat A,\hat B] = c_i^\dagger\big(a_{ik}b_{kj}-b_{ik}a_{kj}\big)c_j
= c_i^\dagger\big([\hat a,\hat b]\big)_{ij}c_j ,
\qquad \langle i|\hat a|j\rangle \equiv a_{ij},
\end{equation}
so the commutator of two bilinears is the bilinear of the \emph{first-quantized}
commutator. The many-body problem in the $2^N$-dimensional Fock space therefore
collapses onto $N\times N$ first-quantized operators. Re-exponentiating the nested
commutators,
\begin{equation}
\label{eq:iS}
iS = \int \mathrm{d} t\;
\big\langle \mathrm{FS}\big|c_i^\dagger c_j\big|\mathrm{FS}\big\rangle\,
\big[\hat U^{-1}\big(-\partial_t - i\hat\epsilon\big)\hat U\big]_{ij},
\qquad \hat U \equiv e^{i\hat\phi},
\end{equation}
with $\hat\epsilon$ the one-particle Hamiltonian
($\hat H = c_i^\dagger \epsilon_{ij}c_j$). Introduce the one-particle density
matrix (statistical matrix) of the ground state,
\begin{equation}
\label{eq:f0matrix}
f_{0,ji} = \big\langle \mathrm{FS}\big|c_i^\dagger c_j\big|\mathrm{FS}\big\rangle,
\end{equation}
so that $E_0=\langle\mathrm{FS}|\hat H|\mathrm{FS}\rangle=\operatorname{Tr}(\hat f_0\hat\epsilon)$.
The action becomes a single first-quantized trace,
\begin{equation}
\label{eq:Strace}
S = \int \mathrm{d} t\;\operatorname{Tr}\!\big[\hat f_0\,\hat U^{-1}\big(i\partial_t-\hat\epsilon\big)\hat U\big]
= \int \mathrm{d} t\;\operatorname{Tr}\!\big[\hat f_0\,\hat U^{-1} i\partial_t \hat U - \hat f\,\hat\epsilon\big],
\qquad \hat f \equiv \hat U \hat f_0 \hat U^{-1},
\end{equation}
where the second form uses cyclicity,
$\operatorname{Tr}[\hat f_0\hat U^{-1}\hat\epsilon\hat U]=\operatorname{Tr}[\hat f\,\hat\epsilon]$, and
$\hat f$ (elements $\langle\Psi|c_i^\dagger c_j|\Psi\rangle$) is the one-particle
density matrix of the coherent state.

\subsection{Wigner transform to the phase-space action}

Map first-quantized operators to phase-space symbols through the Wigner transform
\begin{equation}
\label{eq:wigner}
\omega(\bm{x},\bm{p}) = \int \mathrm{d}^dy\; e^{i\bm{p}\cdot\bm y/\hbar}\,
\big\langle \bm{x}-\tfrac{\bm y}{2}\big|\,\hat\omega\,\big|\bm{x}+\tfrac{\bm y}{2}\big\rangle .
\end{equation}
Two properties close the derivation. A trace becomes a phase-space integral,
\begin{equation}
\label{eq:trwigner}
\operatorname{Tr}(\hat a\hat b) = \int \frac{\mathrm{d}^dx\,\mathrm{d}^dp}{(2\pi\hbar)^d}\, a(\bm{x},\bm{p})\,b(\bm{x},\bm{p}),
\end{equation}
and the symbol of a product is the star product of symbols (so the symbol of a
commutator is $i$ times the Moyal bracket, $\eqref{eq:moyal}$). Applying these to
\eqref{eq:Strace},
\begin{equation}
\label{eq:Sfinal}
\boxed{\;
S = \int \mathrm{d} t\int \frac{\mathrm{d}^dx\,\mathrm{d}^dp}{(2\pi\hbar)^d}\,
\Big[\, f_0(\bm{x},\bm{p})\,U^{-1}\!\star i\partial_t U \;-\; f(\bm{x},\bm{p})\,\epsilon(\bm{p})\,\Big],
\qquad f = U\star f_0\star U^{-1}\; }
\end{equation}
where $U(\bm{x},\bm{p})=\exp[i\phi(\bm{x},\bm{p})]$ is the Wigner symbol of $\hat U$. The first
term is $\langle f_0,\,U^{-1}\!\star i\partial_t U\rangle = S_{\mathrm{WZW}}$ and the
second is $-\langle f,\epsilon\rangle = S_H$, so \eqref{eq:Sfinal} is exactly the
target action \eqref{eq:Saction}, with $f$ on the coadjoint orbit \eqref{eq:orbit}.
This completes the derivation of the bosonized action from the free-fermion theory.

\begin{remark}
The two sections meet at the ground state. The first-quantized density matrix
$\hat f_0$ of \eqref{eq:f0matrix} is the equal-time Fermi-sea projector $\hat n$ of
\S\ref{sec:bridge}; its Wigner symbol is $f_0(\bm{x},\bm{p})=\Theta(p_F-|\bm{p}|)$, and the
dispersion $\epsilon(\bm{p})$ here is the $\xi_{\bm{k}}$ of \S\ref{sec:prelim}. The
$G$--$\Sigma$ fluctuations of \S\ref{sec:bridge} and the orbit coordinate
$f=U\star f_0\star U^{-1}$ describe the same particle--hole physics from the two
coherent-state representations; exhibiting their equality is the bridge the
project constructs.
\end{remark}

\section*{Caveats}

\begin{itemize}
\item \textbf{Convention dependence.} The signs of $\Sigma$ and of the
$\operatorname{Tr}(\Sigma G)$ term in \eqref{eq:gsigma}, and the explicit minus in $G_0^{-1}$
\eqref{eq:G0inv}, are tied to the choices \eqref{eq:Gdef} and \eqref{eq:resolution}.
Different references distribute these signs differently; only the Dyson equation
and the saddle \eqref{eq:freesaddle} are invariant. Any comparison with
\S\ref{sec:coadjoint} or with the literature must first align these conventions.
\item \textbf{Regularization of $\operatorname{Tr}\ln$.} The functional determinant is
UV-divergent; only the difference $\operatorname{Tr}\ln(G_0^{-1}+\Sigma)-\operatorname{Tr}\ln G_0^{-1}$ is
finite, and equal-time objects such as the density require a point-splitting or
frequency-sum prescription (the $e^{i\omega_n 0^+}$ convergence factor). These are
left implicit above.
\item \textbf{Measure and contour.} The bosonic measure $\mathcal{D}G\,\mathcal{D}\Sigma$
and the imaginary contour for $\Sigma$ in \eqref{eq:resolution} are stated
formally; a careful treatment (e.g.\ of the Jacobian from $\psi,\bar{\psi}$ to $G$) is
not needed for the saddle structure but matters for fluctuation determinants.
\item \textbf{The two sections use different coherent states.} \S\ref{sec:gsigma}
is built on the Grassmann coherent state (yielding the bilocal $G$--$\Sigma$
action); \S\ref{sec:coadjoint} on the group coherent state of the Fermi surface
(yielding the coadjoint-orbit action). Both are exact rewritings of the same
free-fermion theory, but their equality is \emph{not} demonstrated here --- it is
the bridge the project constructs.
\item \textbf{Real versus imaginary time.} \S\ref{sec:gsigma} is formulated in
imaginary time (Matsubara), \S\ref{sec:coadjoint} in real time $t$, following the
respective sources. Aligning the two --- including the $i\partial_t$ of the WZW
term against the $\partial_\tau$ of the $\operatorname{Tr}\ln$ --- is part of the reconciliation,
not assumed here.
\item \textbf{Moyal versus Poisson.} Section~\ref{sec:coadjoint} keeps the full
Moyal algebra; the Poisson-bracket form of the coadjoint action (e.g.\ in
Delacr\'etaz--Else--Senthil) is its semiclassical truncation
$\hbar\,\bm{\nabla}_{\bm{x}}\!\cdot\!\bm{\nabla}_{\bm{p}}\ll1$. The distinction is physical, not
cosmetic: it is what allows the magnetic extension of the source paper to capture
Landau-level degeneracy.
\item \textbf{Normalization and total-derivative subtleties of the WZW term.} The
overall normalization of $S_{\mathrm{WZW}}$, and a term linear in $\phi$ that is a
total derivative in the infinite plane but becomes a topological $\theta$-term once
the momentum-space directions are compact, are treated in the body of the source
paper (the weak-field theory), not in the zero-field transcription here.
\item \textbf{Regularization of the traces.} Both the $\operatorname{Tr}\ln$ of
\S\ref{sec:gsigma} and the first-quantized $\operatorname{Tr}$ of \S\ref{sec:coadjoint} are over
infinite single-particle Hilbert spaces and require the usual normal-ordering /
point-splitting prescriptions for equal-time objects; these are left implicit.
\end{itemize}

\end{document}
